\chapter{Родитель отрицательного beta-знака асимптотически-свободной моды}
\label{ch:v10-discrete-resolution-af-beta-sign-parent}

Используем соглашение
\begin{equation}
 \beta_g=\frac{b}{16\pi^2}g^3+O(g^5).
\end{equation}
Для неабелевой калибровочной теории с $n_D$ дираковскими и $n_s$
комплексными фундаментальными мультиплетами
\begin{equation}
 b=-\frac{11}{3}C_A+\frac{2}{3}n_D+\frac{1}{6}n_s.
 \label{eq:v10-af-beta-content}
\end{equation}
Отрицательный вклад самодействия калибровочного поля отсутствует у
унаследованного relative-$U(1)$, где строго получено $b=+2$.

Условие требуемого $b=-2$ эквивалентно
\begin{equation}
 -22C_A+4n_D+n_s=-12.
 \label{eq:v10-af-beta-diophantine}
\end{equation}
Оно имеет безаномальные целые реализации. Например,
\begin{equation}
 SU(2):\ (C_A,n_D,n_s)=(2,8,0),\qquad
 SU(3):\ (C_A,n_D,n_s)=(3,13,2).
\end{equation}
Обе дают точно $b=-2$. Напротив, совпадение для $SU(3)$ с $27$
левыми вейлевскими фундаменталами имеет ненулевую кубическую аномалию;
число $27$ из K43/BV-следа не типизировано как такая кратность.

Однако beta-условие не выбирает носитель. Его линейная карта
\begin{equation}
 B_\beta=(-22,4,1)
\end{equation}
имеет ранг $1$ и ядро размерности $2$. Родитель одной невязки
\begin{equation}
 \mathcal P_\beta
 =\frac12(-22C_A+4n_D+n_s+12)^2
\end{equation}
имеет гессиан ранга $1$ с тем же двухмерным ядром. Следовательно, точный
коэффициент не определяет ни алгебру, ни представления, ни кратности.

Девять внутренних и условных кандидатов проверены по шести критериям.
Матрица имеет ранг $6$, оценки
\begin{equation}
 (4,4,4,4,3,4,1,3,5)
\end{equation}
и ни одного полного прохода. Формальная запись $b=-2$ получает $5/6$,
но не имеет независимого родителя.

\begin{theorem}[Отрицательный знак допустим, но не выведен]
Безаномальный неабелев сектор может точно реализовать $b=-2$. Текущий
relative-$U(1)$ этого не делает, а единственное beta-условие оставляет две
плоские моды состава полей. Поэтому физический AF-носитель не выбран.
\end{theorem}

\begin{nogocriterion}[Числовое совпадение кратности не является типизацией]
Нельзя переименовать число $27$ из спектрального или BV-следа в двадцать
семь хиральных цветных мультиплетов. Требуются калибровочная алгебра,
представления, компенсация аномалий и общий родитель кратностей.
\end{nogocriterion}

\begin{protocol}\begin{enumerate}
\item точные условные безаномальные реализации: \textbf{$2/2$};
\item ранг/ядро beta-родителя: \textbf{$1/2$};
\item полные проходы кандидатов: \textbf{$0/9$};
\item физическое происхождение AF-носителя: \textbf{$0/4$};
\item следующий гейт: \textbf{аудит безаномального AF-носителя}.
\end{enumerate}\end{protocol}

Машинный сертификат сохранён в
\path{s2t_v10_cell_birth_four_volume_nonequilibrium_bath_discrete_resolution_transmuted_mode_asymptotically_free_beta_sign_parent_origin_gate_results.json}.